% templates/carbonio_solution.tex
% Chapitre "Solution Zextras Carbonio" — FIXE, pas de variable Jinja2 ici.
% Présentation générale reprise telle quelle du Générateur de DAT.
% Version "légère" : pas de section Version déployée/Licence (non
% pertinent en phase de dimensionnement pré-vente).

\chapter{Solution Zextras Carbonio}

\section{Présentation générale}
Zextras Carbonio est une plateforme collaborative unifiée conçue pour
centraliser l'ensemble des outils de communication et de collaboration
au sein d'une seule solution. Elle intègre les services de messagerie
électronique, calendriers, contacts, gestion des tâches, partage de
fichiers, messagerie instantanée, visioconférence et collaboration
documentaire, accessibles depuis une interface web moderne ainsi que
depuis des applications mobiles.

Pensée autour des principes de souveraineté numérique, la plateforme
permet aux organisations de conserver la maîtrise de leurs données, de
leur infrastructure et de leurs politiques de sécurité, qu'elles
choisissent un déploiement sur site ou dans un environnement cloud
maîtrisé. Son architecture modulaire offre une grande souplesse de
déploiement et d'évolution afin de répondre aux besoins d'organisations
de toutes tailles.

La plateforme met l'accent sur la sécurité des données grâce à des
mécanismes d'authentification avancés, au chiffrement des
communications et à une gestion fine des droits d'accès. Elle s'appuie
sur des standards ouverts afin de favoriser l'interopérabilité avec les
outils et infrastructures existants.

Les fonctionnalités de partage, de coédition et de communication en
temps réel facilitent le travail collaboratif, aussi bien en présentiel
qu'à distance. Les outils d'administration permettent de gérer les
utilisateurs, les domaines, les politiques de sécurité et les
ressources de manière centralisée. La solution intègre également des
mécanismes de sauvegarde, d'archivage et de restauration afin d'assurer
la protection et la continuité des données.

Grâce à son interface intuitive et à son approche centrée sur
l'expérience utilisateur, Zextras Carbonio constitue un environnement
de travail numérique complet, sécurisé, évolutif et résolument orienté
vers la maîtrise des données et la souveraineté des systèmes
d'information.

Carbonio se déploie aussi bien on-premise que chez un hébergeur
partenaire, ce qui permet de maîtriser la localisation des données ---
point d'attention fréquent des clients publics ou soumis à des
contraintes réglementaires (RGPD, hébergement en France/UE).
